\documentclass[11pt]{article}
\usepackage[T1]{fontenc}
\usepackage{lmodern,amsmath,amssymb,amsthm}
\usepackage[margin=1in]{geometry}
\usepackage[colorlinks=true,allcolors=blue]{hyperref}
\newcommand{\sn}{\operatorname{sn}}
\newcommand{\cn}{\operatorname{cn}}
\newcommand{\dn}{\operatorname{dn}}
\newcommand{\Div}{\operatorname{div}}
\newtheorem{theorem}{Theorem}
\newtheorem{lemma}{Lemma}
\newtheorem{corollary}{Corollary}
\title{The least-period criterion for outer focal products\\in elliptic billiards}
\author{Alec Kriebel\\\small Independent researcher\\
\small ORCID: \href{https://orcid.org/0009-0001-9320-500X}{0009-0001-9320-500X}}
\date{16 October 2026}
\begin{document}
\maketitle
\begin{abstract}
Fix a noncircular ellipse and a proper nested confocal elliptical caustic
supporting periodic billiards of least period $N\geq3$. Intersect the boundary
tangents at consecutive impact points. The product of the ordinary distances
from these outer vertices to either original focus is independent of phase
if and only if $4\mid N$. The two constant products then agree. The positive
direction is the primitive interpretation of the invariant $k_{115}$ and is
also proved in Ferudun's contemporaneous note, which we explicitly credit.
A paired-distance divisor proves both directions, including odd periods and
the four-period degree drop. Repeated traversals cannot restore constancy.
Historical priority and independent discovery are not certified.
\end{abstract}

\section{Statement and provenance}
Write
\[
 E:\frac{x^2}{a^2}+\frac{y^2}{b^2}=1,\qquad
 C:\frac{x^2}{\alpha^2}+\frac{y^2}{\beta^2}=1,
 \qquad a>b>0,
\]
where
\begin{equation}\label{pair}
 a^2-\alpha^2=b^2-\beta^2=\lambda,\quad 0<\lambda<b^2,
 \qquad c^2=a^2-b^2=\alpha^2-\beta^2.
\end{equation}
The foci used throughout are $F_\pm=(\pm c,0)$ of $E$ and $C$.
For a billiard with successive impacts $P_j$ on $E$, let $Q_j$ be the
intersection of the tangents to $E$ at $P_j,P_{j+1}$. Assume that the
fixed pair supports a periodic family of \emph{least} period $N\geq3$.
For either focus put
\[
 R_\pm=\prod_{j=0}^{N-1}|Q_j-F_\pm|.
\]
Distances are ordinary positive Euclidean lengths. Phase varies within the
family for this fixed pair and winding.

\begin{theorem}\label{main}
Under \eqref{pair}, each of $R_+$ and $R_-$ is independent of phase if
and only if $4\mid N$. When $4\mid N$, their constants are equal.
Every admissible coprime star winding is included.
\end{theorem}

Reznik, Garcia and Koiller list the positive assertion as $k_{115}$ in
Table~2, p.~5 of their 2020 preprint~\cite{RGKpre}, and Table~2,
p.~345 of the distinct 2021 journal paper~\cite{RGK}. Those editions
give a proof question mark. This is dated provenance, not evidence of
present worldwide openness. Their table does not expressly settle a
repeated-traversal convention. Theorem~\ref{main} uses least period;
an unrestricted interpretation based only on the length of a repeated
list is false, as Section~\ref{repeat} explains.

Ferudun~\cite{Ferudun}, Theorem~1, proves the same positive assertion
with an explicit constant, for the same primitive scope and all coprime
windings. His Jacobi factorization and cyclic-divisor argument closely
overlap the mechanism below. We present an attributed separate treatment
and the necessary-and-sufficient criterion, without claiming firstness or
historical novelty of that criterion. The dated records and limits of
priority are stated in Section~\ref{credit}.

\section{Geometry and the paired distance}
All Jacobi functions below have modulus $k$, not parameter $k$. Put
\[
 k=\frac c\alpha\in(0,1),\qquad k'=\frac\beta\alpha=\sqrt{1-k^2},
 \qquad K=K(k),\quad K'=K(k').
\]
Stachel's canonical parametrization~\cite[Theorem~4.3 and (4.9)]{Stachel}
gives $v\in(0,K)$ such that
\begin{equation}\label{canonical}
 P(u)=(-a\sn u,b\cn u),\qquad
 a=\frac{\alpha\dn v}{\cn v},\quad b=\frac\beta{\cn v},
 \qquad\delta=2v.
\end{equation}
The billiard map is $u\mapsto u+\delta$. After choosing orientation,
least period is equivalent to
\begin{equation}\label{closure}
 \delta=\frac{4K\tau}{N},\qquad
 \gcd(\tau,N)=1,\qquad 0<\tau<N/2.
\end{equation}
The real parametrization has period $4K$; reversal only permutes the
distance factors. For example, the chord through $P(u-v),P(u+v)$ is
\[
 -\frac{x\sn u}{\alpha}+\frac{y\cn u}{\beta}=1,
\]
the tangent to $C$ at $(-\alpha\sn u,\beta\cn u)$.
The billiard dynamics input in \eqref{canonical} is established
background, rather than a new parametrization.

\begin{lemma}\label{geometry}
The boundary tangents at $P(u-v),P(u+v)$ meet at
\begin{equation}\label{Q}
 Q(u)=(-A\sn u,B\cn u),\qquad
 A=\frac{a\dn v}{\cn v}=\frac{a^2}{\alpha},\quad
 B=\frac b{\cn v}=\frac{b^2}{\beta}.
\end{equation}
For real $u$ this point is finite and outside $E$.
\end{lemma}
\begin{proof}
Put $s=\sn u$, $z=\cn u$, $d=\dn u$, $h=\sn v$,
$C_v^0=\cn v$, $D_v^0=\dn v$, and $D_0=1-k^2s^2h^2$.
The tangents have equations
$-X\sn(u\pm v)/a+Y\cn(u\pm v)/b=1$.
The Jacobi addition formulas~\cite[Section~22.8]{DLMF} reduce their
sum and difference to
\[
 -\frac{XsC_v^0D_v^0}{a}+\frac{YzC_v^0}{b}=D_0,
 \qquad -\frac{Xz}{a}-\frac{YsD_v^0}{b}=0.
\]
Their solution is \eqref{Q}. The original determinant, up to sign, is
$2hC_v^0d/(abD_0)>0$. Finally
\[
 \frac{Q_x^2}{a^2}+\frac{Q_y^2}{b^2}-1
 =\frac{h^2\dn^2u}{\cn^2v}>0.
\]
Thus neither focus is an outer vertex.
\end{proof}

The actual outer sequence is $Q(u_0+v+j\delta)$; its fixed offset is
absorbed into phase. In particular $Q(u+2K)=-Q(u)$.
The conic locus $X^2/A^2+Y^2/B^2=1$ is also the polar reciprocal
of $C$ in $E$: its equation is
$\alpha^2X^2/a^4+\beta^2Y^2/b^4=1$.
This construction is classical polarity~\cite[Section~4]{Fouche}.
The locus need not be confocal with $E$; we never replace $F_\pm$
by its own foci.

Define a meromorphic function without choosing complex square roots:
\begin{equation}\label{Hdef}
 H(u)=\big((Q_x-c)^2+Q_y^2\big)
       \big((Q_x+c)^2+Q_y^2\big).
\end{equation}
For real $u$, $H(u)=(|Q(u)-F_+|\,|Q(u)-F_-|)^2>0$.
\begin{lemma}\label{factor}
For a positive constant $C_v$,
\begin{equation}\label{factorH}
 H(u)=C_v\dn^2u\,
       [\dn^2\delta-k^2\cn^2\delta\,\sn^2u].
\end{equation}
This identity holds before imposing any parity condition on $N$.
\end{lemma}
\begin{proof}
Put $\kappa=k^2$, $t=\sn^2v$, $q=1-t$, and
\[
 U=1-2\kappa t+\kappa t^2,\quad
 V=1-2t+\kappa t^2,\quad W=1-\kappa t^2.
\]
Equation~\eqref{Q} gives
\[
 A=\frac{\alpha(1-\kappa t)}q,\quad
 B=\frac{\alpha k'}q,\quad
 A^2-B^2=\frac{\alpha^2\kappa V}{q^2},\quad
 B^2+c^2=\frac{\alpha^2U}{q^2}.
\]
For $S=\sn^2u$, expand \eqref{Hdef} and use
$U+V=2q(1-\kappa t)$ to obtain
\begin{align*}
 H(u)&=\frac{\alpha^4}{q^4}
 [(U+\kappa VS)^2-4\kappa(1-\kappa t)^2q^2S]\\
 &=\frac{\alpha^4}{q^4}(1-\kappa S)(U^2-\kappa V^2S).
\end{align*}
Now $\cn\delta=V/W$ and $\dn\delta=U/W$ by the double-angle
formulas. This proves \eqref{factorH} with $C_v=\alpha^4W^2/q^4>0$.
The paired factorization is also obtained by multiplying Ferudun's
two one-focus factors~\cite[(5)]{Ferudun}.
\end{proof}

\section{The divisor and the exact period criterion}
Work on the compact connected torus
\[
 T=\mathbb C/(2K\mathbb Z+2iK'\mathbb Z),\qquad p=iK'.
\]
The functions $\sn^2u$ and $\dn^2u$ are meromorphic on $T$.
The first has one double pole at $p$; the second has that pole and
a double zero at $p+K$. We use the standard shift
\begin{equation}\label{shift}
 \sn(z+K+iK')=\frac{\dn z}{k\cn z}
\end{equation}
and the periods in~\cite[Tables~22.4.1--3]{DLMF}.
In particular $\dn(u+2K)=\dn u$.

If $\delta\ne K$, the bracket in \eqref{factorH} has one double pole
at $p$ and zeros at $p+K\pm\delta$, by \eqref{shift}.
These zeros are distinct modulo the lattice: their coincidence would
force $2\delta\in2K\mathbb Z$, hence $\delta=K$ since
$0<\delta<2K$. Neither is $p$ or $p+K$ in this case.
The two distinct zeros exhaust its degree-two zero divisor, and so are
simple. Consequently
\begin{equation}\label{divH}
 \Div H=2[p+K]+[p+K+\delta]+[p+K-\delta]-4[p].
\end{equation}
There is exactly one exceptional primitive case. If $\delta=K$,
\eqref{closure} and coprimality force $N=4$, $\tau=1$; the bracket
is the nonzero constant $k'^2$, and instead
\begin{equation}\label{exception}
 \Div H=2[p+K]-2[p].
\end{equation}
This degree drop is essential.

Let $G=\langle\delta\rangle\subset T$. Its order is
\[
 L=\frac N{\gcd(N,2)},
\]
because $\delta/(2K)=2\tau/N$. Define
\[
 \mathcal F(u)=\prod_{j=0}^{L-1}H(u+j\delta),\qquad
 O_x=\sum_{g\in G}[x+g].
\]
For $N\ne4$, summing translated copies of \eqref{divH} gives
\begin{equation}\label{normdiv}
 \Div\mathcal F=4O_{p+K}-4O_p.
\end{equation}
All multiplicities and collisions between different factors are included:
translation by $\delta$ or $-\delta$ permutes $G$.
The two cosets in \eqref{normdiv} coincide if $K\in G$ and are
otherwise disjoint. On the real period circle, $K$ is the unique point
of order two, so
\begin{equation}\label{parity}
 K\in G\quad\Longleftrightarrow\quad L\text{ is even}
 \quad\Longleftrightarrow\quad4\mid N.
\end{equation}
Thus $\mathcal F$ is constant for $4\mid N$: its zero divisor makes it
holomorphic on the compact torus. At $N=4$, \eqref{exception} yields
the same cancellation with coefficient two. Its constant is finite and
positive, since every factor is positive for real $u$.
For the other parities, \eqref{normdiv} is a nonzero divisor, hence
$\mathcal F$ is nonconstant.

\begin{proof}[Proof of Theorem~\ref{main}]
Write $R_\pm(u)=\prod_{j=0}^{N-1}|Q(u+j\delta)-F_\pm|$.
For all $N$, the antipodal \emph{phase shift} gives
\begin{equation}\label{swap}
 R_-(u)=R_+(u+2K),\qquad
 R_+(u)^2R_-(u)^2=\mathcal F(u)^{N/L}.
\end{equation}
If either focal product were constant, both would be constant by the
first identity. The second would make the meromorphic function
$\mathcal F^{N/L}$ constant on a real interval. It has no real poles;
the identity theorem then makes it constant on $T$. For $4\nmid N$
this contradicts its nonzero divisor \eqref{normdiv}.
This proves necessity also for odd $N$, without assuming that an
antipodal point belongs to the same odd primitive orbit.

If $4\mid N$, then $L=N/2$ and coprimality makes $\tau$ odd.
Now $L\delta=2K\tau$ really does pair the vertices inside the orbit:
$Q(u+(j+L)\delta)=-Q(u+j\delta)$. Therefore
\[
 R_+(u)=R_-(u),\qquad R_+(u)^2=\mathcal F(u).
\]
The unique positive square root of the positive constant proves
sufficiency and equality of the two constants. No complex
square-root branch is used.
\end{proof}

The use of compact-torus pole cancellation is classical; compare
Akopyan, Schwartz and Tabachnikov~\cite{AST}. The period condition
depends on the subgroup, not the chosen coprime generator, so includes
every star winding in \eqref{closure}.

\section{Repeated traversals and an exact boundary example}\label{repeat}
\begin{corollary}
For any positive integer $r$, the product on $r$ traversals is
phase independent exactly when the primitive product is.
\end{corollary}
\begin{proof}
The repeated product is $R_\pm(u)^r$ and $x\mapsto x^r$ is injective
on $(0,\infty)$.
\end{proof}

For an explicit same-caustic check, take $a=3$, $b=1$ and
$\lambda=9/16$. Two primitive six-periodic orbits have vertices
\[
 (3,0),\ (9/4,\sqrt7/4),\ (-9/4,\sqrt7/4)
 \quad\text{and their negatives,}
\]
and
\[
 (0,1),\ (-3\sqrt{15}/4,1/4),\ (-3\sqrt{15}/4,-1/4)
 \quad\text{and their negatives,}
\]
respectively, in the indicated order followed by the negatives in order.
The caustic has squared axes $135/16,7/16$.
Boundary membership, tangency and the reflection law can be checked
exactly; the verification package supplies these checks.
The first three outer vertices are, respectively,
\[
 (3,1/\sqrt7),\ (0,4/\sqrt7),\ (-3,1/\sqrt7),
\]
and
\[
 (-9/\sqrt{15},1),\ (-12/\sqrt{15},0),\ (-9/\sqrt{15},-1),
\]
followed by their negatives. The products to $F_+=(\sqrt8,0)$ are
\[
 \frac{20736}{343}\qquad\text{and}\qquad\frac{6912}{125},
\]
which are unequal. Traversing each twice produces nominal length
twelve and unequal squared products. Ferudun~\cite[Section~5]{Ferudun}
gives a different repeated-triangle example for the same interpretive
issue; this distinction is explicitly credited.

For a circle, both foci coincide with the center and the outer radius
of a regular star is $a/\cos(\pi\tau/N)$, so every admitted period has
a constant product. Noncircularity is therefore essential to necessity.
Hyperbolic or collapsed caustics, separatrices and $N=2$ are outside
Theorem~\ref{main}; no conclusion about them is drawn by a limit.

\section{Attribution, timeline and verification}\label{credit}
The submitted positive-direction proof is retained at
\href{https://github.com/AlecKriebel/Math/pull/200}{GitHub PR~200},
whose server creation record is 1~October~2026, 05:03:51~UTC, with
submitted head \texttt{379a95960995}.
Ferudun's deposited record~\cite{Ferudun} reports creation on
2~October~2026, 01:41:04.506303~UTC; DataCite records DOI registration
at 01:41:05~UTC that day. These particular witnesses document
contemporaneous separate treatments. They do not exclude an earlier
website or other repository disclosure, establish absolute first
publication, or certify independent discovery. The converse criterion
was added and audited on 16~October, and is not backdated to the original
submission. A bounded classical and modern literature audit did not
establish its historical novelty.

Canonical parametrization, polarity, Jacobi identities and the norm
principle are credited background. Older alternating focal-product
relations, such as~\cite[Section~2, p.~219]{Anonymous}, compare the two
foci for a fixed circumscribed polygon; that assertion alone is not
phase constancy of the full product. The supplementary priority ledger
records exact source scopes and access/search limits. No exhaustive
absence-of-prior-work claim is made.

The supplementary verification includes exact coefficient identities,
cyclic divisor arithmetic, physical reflection/tangency fixtures and
high-precision direct tangent-intersection controls. Numerical samples
are diagnostics, not the universal proof. Legacy assertion-based
verifiers and their optimized-mode limitation are distinguished from
the package's explicit-guard checker. Dependency versions and replay
instructions are supplied. No formal proof-assistant certification is
claimed.

AI tools were used extensively in developing the original argument,
extending and checking it, searching sources, drafting the manuscript
and conducting adversarial audits. These audits are AI-assisted checks,
not conventional independent human peer review. This is an unrefereed
research note; the author is responsible for its claims and presentation.

\begin{thebibliography}{9}
\bibitem{RGKpre} D.~Reznik, R.~Garcia and J.~Koiller,
\emph{Eighty New Invariants in the Elliptic Billiard},
\href{https://arxiv.org/abs/2004.12497v11}{arXiv:2004.12497v11} (2020).
\bibitem{RGK} D.~Reznik, R.~Garcia and J.~Koiller,
Fifty New Invariants of $N$-Periodics in the Elliptic Billiard,
\emph{Arnold Math. J.} \textbf{7} (2021), 341--355.
\href{https://doi.org/10.1007/s40598-021-00174-y}{doi:10.1007/s40598-021-00174-y}.
\bibitem{Ferudun} A.~Ferudun,
\emph{Focal Distance Products of Outer Elliptic Billiard Polygons},
unrefereed preprint (2026).
\href{https://doi.org/10.5281/zenodo.23093258}{doi:10.5281/zenodo.23093258}.
\bibitem{Stachel} H.~Stachel, On the motion of billiards in ellipses,
\emph{Eur. J. Math.} \textbf{8} (2022), 1602--1622.
\href{https://doi.org/10.1007/s40879-021-00524-2}{doi:10.1007/s40879-021-00524-2}.
\bibitem{DLMF} NIST Digital Library of Mathematical Functions,
Sections 22.4, 22.6 and 22.8,
\url{https://dlmf.nist.gov/22} (accessed 16 October 2026).
\bibitem{Fouche} M.~Fouch\'e, Sur les polygones inscrits \`a une conique
et circonscrits \`a une autre et les polygones form\'es de g\'en\'eratrices
d'une quadrique et inscrits \`a une autre quadrique,
\emph{Bull. Soc. Math. France} \textbf{43} (1915), 146--160.
\url{https://www.numdam.org/item/BSMF_1915__43__146_1/}.
\bibitem{AST} A.~Akopyan, R.~Schwartz and S.~Tabachnikov,
\emph{Billiards in ellipses revisited},
\href{https://arxiv.org/abs/2001.02934v2}{arXiv:2001.02934v2} (2020).
\bibitem{Anonymous} Anonymous, Th\'eor\`eme sur les polygones
circonscrits \`a une conique, \emph{Nouvelles Annales de Math\'ematiques},
1st series, \textbf{12} (1853), 218--220.
\url{https://www.numdam.org/item/NAM_1853_1_12__218_1/}.
\end{thebibliography}
\end{document}
